\subsection{Nonnumeric Peano Systems}

\begin{topicbox}{Different Carriers, Same Successor Shape}
\begin{exposition}
A Peano carrier need not be the usual natural-number set, and its successor
operation need not agree with the ambient order or arithmetic of a larger
structure. What matters is the internally generated successor chain.
\end{exposition}

\begin{example*}[The Powers of Ten]
Let
\[
P=\{10^{n-1}:n\in\mathbb{N}\}, \qquad e=1,
\]
and define
\[
S:P\to P,
\qquad
S(10^{n-1})=10^n
\quad(n\in\mathbb{N}).
\]
Equivalently, for every $x\in P$,
\[
S(x)=10x.
\]
Then $(P,S,e)$ is a Peano system.
\end{example*}

\begin{remark*}[Structure]
The Peano data are
\[
\bigl(P,S,e\bigr)
=
\bigl(\{1,10,10^2,10^3,\ldots\},\,x\mapsto 10x,\,1\bigr).
\]
\end{remark*}

\begin{remark*}[Interpretation]
The elements are not all natural numbers. They are only powers of ten.
Nevertheless, starting from $1$ and repeatedly multiplying by $10$ produces a
successor chain with exactly the Peano shape:
\[
1 \mapsto 10 \mapsto 10^2 \mapsto 10^3 \mapsto \cdots.
\]
\end{remark*}

\begin{example*}[Strings of Repeated Letters]
Let $\Sigma=\{a\}$ and let
\[
P=\{a^n:n\in\mathbb{N}\}, \qquad e=a,
\]
where $a^n$ denotes the string consisting of $n$ copies of $a$. Define
\[
S:P\to P,
\qquad
S(a^n)=a^{n+1}
\quad(n\in\mathbb{N}).
\]
Then $(P,S,e)$ is a Peano system.
\end{example*}

\begin{remark*}[Structure]
The Peano data are
\[
\bigl(P,S,e\bigr)
=
\bigl(\{a,aa,aaa,aaaa,\ldots\},\,a^n\mapsto a^{n+1},\,a\bigr).
\]
\end{remark*}

\begin{remark*}[Interpretation]
This example is useful because the elements are not numbers at all. They are
finite strings. The Peano structure is carried by the successor relation, not
by the surface appearance of the elements.
\end{remark*}

\begin{example*}[The Reciprocal Chain]
Let
\[
P=\left\{\frac{1}{n}:n\in\mathbb{N}\right\}, \qquad e=1,
\]
and define
\[
S:P\to P,
\qquad
S\!\left(\frac{1}{n}\right)=\frac{1}{n+1}
\quad(n\in\mathbb{N}).
\]
Then $(P,S,e)$ is a Peano system.
\end{example*}

\begin{remark*}[Structure]
The Peano data are
\[
\bigl(P,S,e\bigr)
=
\left(
\left\{1,\frac{1}{2},\frac{1}{3},\frac{1}{4},\ldots\right\},
\,\frac{1}{n}\mapsto \frac{1}{n+1},
\,1
\right).
\]
\end{remark*}

\begin{remark*}[Interpretation]
This is a good example precisely because the carrier is a decreasing sequence
in the usual order on $\mathbb{R}$:
\[
1>\frac{1}{2}>\frac{1}{3}>\frac{1}{4}>\cdots.
\]
The Peano successor does not have to mean ``larger than.'' It only means the
next element in the generated successor chain. The order inherited from
$\mathbb{R}$ is external structure, not part of the Peano-system data.
\end{remark*}
\end{topicbox}

